\section{Resultados e Discussão}

\begin{frame}{Métricas médias em 100 execuções}
  \centering
  \begin{adjustbox}{max width=\textwidth}
\scriptsize
\setlength{\tabcolsep}{4pt}
\renewcommand{\arraystretch}{1.15}
% Mean followed by its standard deviation, the latter in a slightly smaller font (7 pt vs. 8 pt)
\newcommand{\meanstd}[2]{#1\,{\fontsize{7}{8}\selectfont$\pm$\,#2}}
\begin{tabular}{>{\cellcolor{white}}c l c c c c c c}
    \toprule
    & &
    \multicolumn{4}{c}{Resposta do punho} &
    \multicolumn{2}{c}{Sinal de controle} \\
    \cmidrule(lr){3-6} \cmidrule(lr){7-8}
    & Controlador &
    TPSR $\uparrow$ & $R^2$ $\uparrow$ & Entropia $\downarrow$ & IAE $\downarrow$ & Potência $\downarrow$ & TV $\downarrow$ \\
    & &
    [\%] & [---] & [bits] & [rad$\cdot$s] & [(N$\cdot$m)$^2$] & [N$\cdot$m$\cdot$s] \\
    \midrule
    & EADRC+EBMFLC & \meanstd{45,92}{54,00} & \meanstd{0,41}{0,15} & \meanstd{5,18}{0,16} & \meanstd{0,83}{0,28} & \meanstd{0,50}{0,26} & \meanstd{307,44}{89,98} \\
    \rowcolor{CBALightPurple}
    & EADRC+ZPLP & \meanstd{57,49}{35,50} & \meanstd{0,46}{0,20} & \meanstd{5,17}{0,19} & \meanstd{0,76}{0,20} & \meanstd{0,34}{0,12} & \meanstd{263,51}{47,82} \\
    & DE-PID & \meanstd{99,90}{0,10} & \meanstd{1,00}{0,00} & \meanstd{4,85}{0,10} & \meanstd{0,03}{0,01} & \meanstd{0,98}{0,38} & \meanstd{877,68}{198,91} \\
    \multirow{-4}{*}{\rotatebox{90}{Repouso}}
    & IMC-PID & \meanstd{$-$123,32}{563,17} & \meanstd{$-$0,37}{1,85} & \meanstd{5,50}{0,31} & \meanstd{1,24}{1,03} & \meanstd{8,75}{25,33} & \meanstd{492,03}{260,60} \\
    \midrule
    & EADRC+EBMFLC & \meanstd{45,02}{54,72} & \meanstd{0,46}{0,12} & \meanstd{5,23}{0,13} & \meanstd{0,84}{0,28} & \meanstd{0,51}{0,26} & \meanstd{308,50}{89,95} \\
    \rowcolor{CBALightPurple}
    & EADRC+ZPLP & \meanstd{56,97}{35,77} & \meanstd{0,51}{0,16} & \meanstd{5,21}{0,15} & \meanstd{0,76}{0,20} & \meanstd{0,35}{0,12} & \meanstd{264,31}{47,75} \\
    & DE-PID & \meanstd{99,90}{0,10} & \meanstd{1,00}{0,00} & \meanstd{4,92}{0,09} & \meanstd{0,03}{0,01} & \meanstd{0,98}{0,38} & \meanstd{877,68}{198,91} \\
    \multirow{-4}{*}{\rotatebox{90}{Movimento}}
    & IMC-PID & \meanstd{$-$154,62}{684,56} & \meanstd{$-$0,40}{2,17} & \meanstd{5,50}{0,29} & \meanstd{1,31}{1,12} & \meanstd{10,39}{31,63} & \meanstd{511,02}{295,24} \\
    \bottomrule
\end{tabular}
\end{adjustbox}


  {\scriptsize Média $\pm$ desvio padrão; $\uparrow$ ($\downarrow$): maior (menor) é melhor.}

  \begin{itemize}[<+->]
    \small
    \item DE-PID: TPSR $\approx$ 99,9\%, mas conhece $\theta_{v_3}$ e tem TV $\approx 3{,}3\times$ a do EADRC+ZPLP.
    \item EADRC+ZPLP: maior TPSR entre os controladores realizáveis (56,97\% em movimento), com menor esforço que o EADRC+EBMFLC.
    \item IMC-PID: TPSR mediano acima de 50\%, mas oscilações crescentes em algumas execuções derrubam a média.
  \end{itemize}
\end{frame}

\begin{frame}{Análise estatística}
  \begin{columns}[c]
    \begin{column}{0.62\textwidth}
      \centering
      \scriptsize
\setlength{\tabcolsep}{4pt}
\begin{tabular}{c l cc cc cc}
    \toprule
    & & \multicolumn{6}{c}{EADRC+ZPLP vs.\ $\cdots$} \\
    \cmidrule(lr){3-8}
    & & \multicolumn{2}{c}{EADRC+EBMFLC} & \multicolumn{2}{c}{DE-PID} & \multicolumn{2}{c}{IMC-PID} \\
    \cmidrule(lr){3-4} \cmidrule(lr){5-6} \cmidrule(lr){7-8}
    & Métrica & $p$ & $r$ & $p$ & $r$ & $p$ & $r$ \\
    \midrule
    & TPSR     & $<$0,001 & +0,56 & $<$0,001 & $-$1,00 & $<$0,001 & +1,00 \\
    & $R^2$    & $<$0,001 & +0,61 & $<$0,001 & $-$1,00 & $<$0,001 & +1,00 \\
    & Entropia & $<$0,001 & $-$0,69 & $<$0,001 & +1,00 & $<$0,001 & $-$1,00 \\
    & IAE      & $<$0,001 & $-$0,67 & $<$0,001 & +1,00 & $<$0,001 & $-$1,00 \\
    & Potência & $<$0,001 & $-$1,00 & $<$0,001 & $-$1,00 & \alert{0,619} & \alert{+0,06} \\
    \multirow{-6}{*}{\rotatebox{90}{Repouso}}
    & TV       & $<$0,001 & $-$0,85 & $<$0,001 & $-$1,00 & $<$0,001 & +1,00 \\
    \midrule
    & TPSR     & $<$0,001 & +0,60 & $<$0,001 & $-$1,00 & $<$0,001 & +1,00 \\
    & $R^2$    & $<$0,001 & +0,67 & $<$0,001 & $-$1,00 & $<$0,001 & +1,00 \\
    & Entropia & $<$0,001 & $-$0,86 & $<$0,001 & +1,00 & $<$0,001 & $-$1,00 \\
    & IAE      & $<$0,001 & $-$0,69 & $<$0,001 & +1,00 & $<$0,001 & $-$1,00 \\
    & Potência & $<$0,001 & $-$1,00 & $<$0,001 & $-$1,00 & $<$0,001 & $-$0,70 \\
    \multirow{-6}{*}{\rotatebox{90}{Movimento}}
    & TV       & $<$0,001 & $-$0,86 & $<$0,001 & $-$1,00 & $<$0,001 & +1,00 \\
    \bottomrule
\end{tabular}


      \vspace{0.1cm}
      {\scriptsize $p$ corrigido por Holm-Bonferroni; $r > 0$: métrica maior com EADRC+ZPLP; em vermelho, $p \geq 0{,}05$.}
    \end{column}
    \begin{column}{0.36\textwidth}
      \small
      \begin{itemize}[<+->]
        \setlength{\itemsep}{0.2cm}
        \item \textbf{vs.\ EBMFLC:} maior supressão (TPSR, IAE) e \textbf{menor esforço} de controle (potência menor em todas as execuções); entropia equivalente em repouso.
        \item \textbf{vs.\ DE-PID:} $|r| \approx 1$ em todas as métricas: resposta melhor em todas as execuções, à custa de maior esforço em todas elas.
        \item \textbf{vs.\ IMC-PID:} supressão superior ($|r| \geq 0{,}94$ em TPSR, $R^2$ e IAE) e \textbf{menor esforço} em todas as execuções (potência e TV).
      \end{itemize}
    \end{column}
  \end{columns}
\end{frame}

\begin{frame}{Espectrogramas da resposta do punho}
  \begin{columns}[c]
    \begin{column}{0.68\textwidth}
      \centering
      \begin{minipage}{0.49\linewidth}
        \centering
        {\scriptsize Sem controle, repouso}\\
        \includegraphics[width=\linewidth]{figures/plots/uncontrolled_amplitude_0.0/spectrogram_response_uncontrolled_amplitude_0.0.pdf}
      \end{minipage}
      \hfill
      \begin{minipage}{0.49\linewidth}
        \centering
        {\scriptsize EADRC+EBMFLC, repouso}\\
        \includegraphics[width=\linewidth]{figures/plots/eadrc_ebmflc_amplitude_0.0/spectrogram_response_eadrc_ebmflc_amplitude_0.0.pdf}
      \end{minipage}

      \vspace{0.2cm}
      \begin{minipage}{0.49\linewidth}
        \centering
        {\scriptsize Sem controle, movimento}\\
        \includegraphics[width=\linewidth]{figures/plots/uncontrolled_amplitude_1.0/spectrogram_response_uncontrolled_amplitude_1.0.pdf}
      \end{minipage}
      \hfill
      \begin{minipage}{0.49\linewidth}
        \centering
        {\scriptsize EADRC+EBMFLC, movimento}\\
        \includegraphics[width=\linewidth]{figures/plots/eadrc_ebmflc_amplitude_1.0/spectrogram_response_eadrc_ebmflc_amplitude_1.0.pdf}
      \end{minipage}

      \vspace{0.1cm}
      {\tiny Espectrogramas de simulações longas (cerca de 1000 s), para maior resolução em frequência.}
    \end{column}
    \begin{column}{0.30\textwidth}
      \small
      \begin{itemize}
        \setlength{\itemsep}{0.2cm}
        \item Componente de 14,35 Hz fortemente atenuada.
        \item Componentes de 3,59 e 5,30 Hz \alert{pouco mitigadas}.
        \item Comportamento semelhante em repouso e em movimento.
      \end{itemize}
    \end{column}
  \end{columns}
\end{frame}

